\section{Introduction}
\label{sec:intro}

Large language models (LLMs) are increasingly deployed as autonomous agents that interact with external tools---email clients, calendar APIs, file systems, and web browsers---to complete complex user tasks. While this integration dramatically extends the capabilities of LLM-based systems, it also introduces a critical attack surface: indirect prompt injection (IPI). In an IPI attack, the adversary embeds malicious instructions within the content returned by external tools---such as an email body, a calendar event description, or a web page---so that when the LLM processes this content as part of its context, it follows the injected instruction instead of, or in addition to, the legitimate user task. Unlike direct prompt injection, where the attacker controls the user's input, IPI exploits the fundamental inability of current LLMs to distinguish instructions from data when both are presented as text in the same context window.

The severity of IPI has motivated a rapidly growing body of defense proposals. These defenses span a wide design space: input-encoding approaches wrap untrusted content in delimiters or encodings~\cite{hines2024defending}; structured-query defenses separate trusted instructions from untrusted data at the prompt level~\cite{chen2024struq}; runtime-checking defenses verify that proposed actions align with the original user task~\cite{jia2025task}; internal-probing defenses inspect the LLM's attention patterns for signs of injection~\cite{hung2025attention}; and architecture-level defenses use process isolation or trusted execution environments to contain the blast radius of a hijacked agent~\cite{wu2025isolategpt,li2025security}. Training-based defenses fine-tune the LLM to resist injection via preference optimization~\cite{chen2024secalign}.

Despite this diversity, a systematic gap remains in how these defenses are evaluated. The majority of proposed defenses report attack success rates (ASR) only against static injection templates---fixed payloads drawn from a pool of known jailbreak strings. Recent work by Zhan et al.~\cite{zhan2025adaptive} demonstrated that adaptive attacks, which tailor the injection payload to the specific defense's detection signal, break three representative defenses (Spotlighting, StruQ, and SecAlign), raising ASR from below 10\% to over 60\%. This finding raises a pressing question: \emph{does the collapse generalize beyond three defenses, and if so, can the failure modes be systematized into root causes and design principles?}

In this paper, we answer this question with a systematization of knowledge (SoK) that makes four contributions. First, we develop a two-dimensional defense taxonomy that classifies 13 existing IPI defenses by their layer position (where in the agent pipeline they intervene) and trust assumption (what they assume they can trust). Second, we build a unified adaptive attack framework that, given a defense class, materializes a matched adversary using one of six evasion strategies---each targeting the specific signal the defense relies on. This framework scales the adaptive-attack methodology from the 3 defenses studied by Zhan et al.\ to all 13 defenses in our taxonomy. Third, we conduct a large-scale evaluation on the AgentDojo benchmark~\cite{debenedetti2024agentdojo}, running 3{,}900 trials (13 defenses $\times$ 5 attack levels $\times$ 3 seeds $\times$ 20 tasks) and find that the failure modes cluster cleanly into four root-cause families, each dominant for a specific defense class. Fourth, we formalize a single-signal impossibility theorem---any defense whose decision is a function of a single signal computable from the prompt can be broken by adaptive search---and instantiate four design-principle prototypes to empirically test the theorem's predictions.

Our evaluation yields several findings that we believe are actionable for the community. Defenses based on input encoding, internal probing, and architecture-level containment all collapse under adaptive attack (ASR rises from 0.73--0.88 to 1.00), and their failures are predominantly attributable to a single root cause: observable-signal mimicry (R1), where the attacker reuses the defense's own delimiter, encoding, or attention pattern. Runtime-checking defenses achieve 0\% adaptive ASR, but only at the cost of blocking 47--100\% of benign actions, rendering them impractical. The only defense class that resists adaptive attack without catastrophic utility loss is structured queries (StruQ, ASR 0.10), but this comes with a 98\% reduction in static ASR, suggesting it may simply be too strong an encoding for the LLM to follow any instruction---benign or malicious. Our design-principle prototypes reveal that a single unobservable signal (P1) is insufficient (R-ASR 0.90), while orthogonal multi-signal combinations (P2) resist but incur a 57\% false-positive rate, quantifying the fundamental security-utility tradeoff.

We release our framework, evaluation code, and the AdaptiveIPI-Bench benchmark as open-source artifacts to support reproducible evaluation of future IPI defenses.

\paragraph{Contributions.}
\begin{itemize}[noitemsep,topsep=0pt]
\item A two-dimensional taxonomy of 13 IPI defenses across 5 defense classes (\Cref{sec:taxonomy}).
\item A unified adaptive attack framework with 6 evasion strategies, scaling Zhan et al.'s 3-defense study to all 13 defenses (\Cref{sec:attack}).
\item A large-scale evaluation (3{,}900 trials) revealing 4 root-cause families that cluster by defense class (\Cref{sec:eval}).
\item A single-signal impossibility theorem and 4 design-principle prototypes (\Cref{sec:principles}).
\end{itemize}
